% 原书物理页=印刷页36--37；1.11 全节.（末尾未编号定理的结论 i)、ii) 及其后两个例子实际印于原书第 38 页顶部、1.12 节标题之前，属本节内容，一并译出附于末尾.）
\section{复 Lie 群. 复化.}\label{sec:lie-1-11}

\begin{numpara}[复 Lie 群]\label{np:lie-1-11-1}
一个 \newterm{lie-complex-lie-group}{复 Lie 群}{complex Lie group} 是一个 \newterm{lie-complex-manifold}{复流形}{complex manifold}，其上赋予了群结构，使得乘法 $m$ 与求逆 $i$ 都是 \newterm{lie-holomorphic}{全纯的}{holomorphic}.
\noindent\textbf{例.}\ $\mathrm{GL}_n(\mathbb C)$、$\mathrm{Sp}_{2n}(\mathbb C)$、$\mathrm{SO}_n(\mathbb C)$，但 $\mathrm{U}(n)$ 不是！
若 $V$ 是一个复流形，那么它的底光滑流形 $M$ 带有一个``\newterm{lie-almost-complex-structure}{殆复结构}{almost complex structure}''，即 $\operatorname{End}_{\mathbb R}(TM)$ 的一个满足 $J^{2}=-\operatorname{id}$ 的截面 $J$（它在每个切空间 $T_xM$ 上给出乘以 $i$ 的运算.） 由此推知，复 Lie 群的 Lie 代数是一个 $\mathbb C$-Lie 代数.
反过来，设 $G$ 是一个 Lie 群，其 Lie 代数带有复结构（从而是一个 $\mathbb C$-Lie 代数），那么可以通过平移把它扩充为 $G$ 上的一个殆复结构.
\noindent\textbf{事实.}\ \emph{这一殆复结构来自 $G$ 上的一个（唯一的）复结构. 此外，复 Lie 群之间的光滑态射 $\varphi:G\longrightarrow G'$ 是全纯的，当且仅当 $d_e\varphi$ 是 $\mathbb C$-线性的.}
第一条论断是殆复结构``可积性''的 Newlander--Nirenberg 判据的一个推论. 第二条则更容易.
\begin{proposition*}
任何紧连通的复 Lie 群都是交换群，且形如 $\mathbb C^{n}/\Gamma$，其中 $\Gamma$ 是 $\mathbb C^{n}$ 的一个余紧 \newterm{lie-lattice}{格}{lattice}.
\end{proposition*}
\begin{proof}
$g\mapsto\operatorname{Ad}_g\in\mathfrak{gl}_n(\operatorname{Lie}(G))$ 的系数都是全纯的，由于 $G$ 紧且连通，故它们是常数. 于是对一切 $g$ 有 $\operatorname{Ad}_g=\operatorname{id}_{\operatorname{Lie}(G)}$，再由连通性，对一切 $g$ 有 $\operatorname{Int}_g=\operatorname{id}_G$. 从而 $G$ 是交换的. 由此可知 $\exp_G:\operatorname{Lie}(G)\longrightarrow G$ 是一个 Lie 群态射，并且由于它在 $0$ 处的微分是恒等映射，故它是一个覆盖. 其核 $\Gamma$ 是离散且余紧的，并且有 $\operatorname{Lie}(G)/\Gamma\xrightarrow{\ \sim\ }G$.
\end{proof}
\end{numpara}

\begin{numpara}[实 Lie 代数的复化]\label{np:lie-1-11-2}
如果 $\mathfrak g$ 是一个 $\mathbb R$-Lie 代数，那么 $\mathfrak g_{\mathbb C}:=\mathbb C\otimes_{\mathbb R}\mathfrak g$ 自然地是一个 $\mathbb C$-Lie 代数，称为 $\mathfrak g$ 的``\newterm{lie-complexification}{复化}{complexification}''. 偶对 $(\mathfrak g_{\mathbb C},\iota:\mathfrak g\hookrightarrow\mathfrak g_{\mathbb C})$ 具有如下泛性质：对任何 $\mathbb C$-Lie 代数 $\mathfrak h$，映射
\[
\begin{array}{rcl}
\operatorname{Hom}_{\mathbb C-AL}(\mathfrak g_{\mathbb C},\mathfrak h) & \longrightarrow & \operatorname{Hom}_{\mathbb R-AL}(\mathfrak g,\mathfrak h)\\
f & \longmapsto & f\circ\iota
\end{array}
\]
是一个双射. 这由张量积的泛性质得出. 我们也说 $\mathfrak g$ 是 $\mathfrak g_{\mathbb C}$ 的一个 \newterm{lie-real-form}{实形式}{real form}.
\begin{example*}
$\mathfrak{gl}_n(\mathbb R)$ 显然是 $\mathfrak{gl}_n(\mathbb C)$ 的一个实形式，但 $\mathfrak u(n)$ 也是它的一个实形式. 事实上，映射
\[
\begin{array}{rcl}
\mathbb C\otimes_{\mathbb R}\mathfrak u(n) & \longrightarrow & \mathfrak{gl}_n(\mathbb C)\\
z\otimes M & \longmapsto & zM
\end{array}
\]
是一个 $\mathbb C$-Lie 代数的同构. 其逆同构按如下方式给出. 若 $M\in\mathrm{M}_n(\mathbb C)$，则把它写成 $M=A+H$ 的形式，其中 $A$ 是反埃尔米特矩阵，$H$ 是埃尔米特矩阵（于是 $A=\frac{1}{2}(M-{}^{t}\bar M)$，$H=\frac{1}{2}(M+{}^{t}\bar M)$），并把 $M$ 送到 $\mathbb C\otimes_{\mathbb R}\mathfrak u(n)$ 中的 $1\otimes A-i\otimes iH$.
\end{example*}
\end{numpara}

\begin{numpara}[Lie 群的复化]\label{np:lie-1-11-3}
对于 Lie 群，没有办法像对 Lie 代数那样容易地进行数量扩张. 尽管如此，前面的讨论提示了如下定义：
\begin{definition*}
设 $G$ 是一个 Lie 群. 由一个复 Lie 群 $G_c$ 与一个 Lie 群态射组成的偶对 $(G_c,\iota:G\longrightarrow G_c)$ 称为 $G$ 的一个复化，如果它具有如下泛性质：对任何复 Lie 群 $H$，映射
\[
\begin{array}{rcl}
\operatorname{Hom}_{GpLieComp}(G_c,H) & \longrightarrow & \operatorname{Hom}_{GpLie}(G,H)\\
\varphi & \longmapsto & \varphi\circ\iota
\end{array}
\]
是双射的.
\end{definition*}
泛性质保证了可能存在的复化的唯一性. 至于存在性，可以引用下面这个我们不予证明的定理：
\begin{theorem*}
任何紧连通 Lie 群 $G$ 都具有一个复化 $(G_c,\iota)$. 此外：
\begin{enumerate}[label=\roman*)]
\item $\iota$ 是一个嵌入，且 $\iota(G)$ 是 $G_c$ 的一个极大紧子群；
\item $(\operatorname{Lie}(G_c),d\iota)$ 是 $\operatorname{Lie}(G)$ 的复化.
\end{enumerate}
我们也说 $G$ 是 $G_c$ 的一个实形式.
\end{theorem*}
\begin{example*}
$\mathrm{SU}(n)$ 与 $\mathrm{SL}_n(\mathbb R)$ 都是 $\mathrm{SL}_n(\mathbb C)$ 的实形式. 对于 $\mathrm{SU}(n)$，这由下面的图得出（其中 $G$ 表示一个复 Lie 群）
\[
\begin{tikzcd}
\operatorname{Hom}_{GpLieComp}(\mathrm{SL}_n(\mathbb C),G)
\ar[r,"\mathrm{res}"]
\ar[d]
&
\operatorname{Hom}_{GpLie}(\mathrm{SU}(n),G)
\ar[d]
\\
\operatorname{Hom}_{\mathbb C-AL}(\mathfrak{sl}_n(\mathbb C),\operatorname{Lie}(G))
\ar[r,"\mathrm{res}"]
&
\operatorname{Hom}_{\mathbb R-Al}(\mathfrak{su}(n),\operatorname{Lie}(G))
\end{tikzcd}
\]
其中下面的箭头是双射的，因为我们已经看到 $\mathfrak{su}(n)$ 是 $\mathfrak{sl}_n(\mathbb C)$ 的一个实形式；而竖直的箭头是双射的，因为 $\mathrm{SU}(n)$ 与 $\mathrm{SL}_n(\mathbb C)$ 都是单连通的. 对于 $\mathrm{SL}_n(\mathbb R)$，可以作出同样的图，但我们不能先验地知道右边的箭头是满射，因为 $\mathrm{SL}_n(\mathbb R)$ 不是单连通的. 不过我们知道它是单射的，这已足以断定上面的箭头是双射的.
\end{example*}
\begin{example*}
下面的嵌入都是复化：
\begin{enumerate}[label=\roman*)]
\item $S^{1}\hookrightarrow\mathbb C^{\times}$；
\item $\mathrm{U}(n)\hookrightarrow\mathrm{GL}_n(\mathbb C)$；
\item $\mathrm{SO}(n)\hookrightarrow\mathrm{SO}(n,\mathbb C)$；
\item $\mathrm{GL}_n(\mathbb R)\hookrightarrow\mathrm{GL}_n(\mathbb C)$.
\end{enumerate}
对于 i)，设 $\varphi:S^{1}\longrightarrow G$ 是到一个复群 $G$ 的态射，令 $X:=\frac{1}{2i\pi}d\varphi(1)$. 那么，态射 $\mathbb C\longrightarrow G$，$z\mapsto\exp_G(zX)$ 是一个全纯的群态射，它在 $2i\pi\mathbb Z$ 上是平凡的. 于是它经由商 $\mathbb C^{\times}=\mathbb C/2i\pi\mathbb Z$（由 $z\mapsto e^{z}$ 所诱导）分解为一个态射 $\mathbb C^{\times}\longrightarrow G$，后者在把 $S^{1}$ 等同于 $\{z\in\mathbb C^{\times},\ |z|=1\}$ 时延拓了 $\varphi$.
\end{example*}
\end{numpara}
